\clearpage
\VerdierPhysicalPageAnchor{0210}%
\begingroup\small\setlength{\parskip}{0pt}\setlength{\emergencystretch}{1.5em}
\setlength{\abovedisplayskip}{4pt plus 2pt minus 2pt}%
\setlength{\belowdisplayskip}{4pt plus 2pt minus 2pt}%

\VerdierSemanticAnchor{ch3d-3-1-4-proof}
只须证明第三项断言。其余断言都是命题 (2.4.4) 的直接推论。断言 2) 立即
表明，$\mathcal A$ 的上同调维数小于或等于 $\mathcal A$ 的对象的内射维数（相应地，
投射维数）的上确界。设 $\mathcal A$ 的上同调维数等于 $n$，并设

\VerdierSemanticAnchor{ch3d-3-1-4-proof-resolution}
\[
I^\bullet=\cdots\longrightarrow0\longrightarrow0\longrightarrow I^0
\longrightarrow I^1\longrightarrow I^2\longrightarrow I^3
\longrightarrow\cdots
\]
是 $\mathcal A$ 的一个对象 $Y$ 的内射型分解。对 $\mathcal A$ 的每个对象
$X$，都有
$\operatorname{Hom}_{K(\mathcal A)}(X[-n-1],I^\bullet)=0$。特别地，取
$X=\operatorname{Ker}d_{I^\bullet}^{n+1}$，立即可知典范态射
$\operatorname{Ker}d_{I^\bullet}^{n+1}\longrightarrow I^{n+1}$ 通过
$I^n$ 分解：

\VerdierSemanticAnchor{ch3d-3-1-4-proof-factorization}
\[
\begin{tikzcd}[column sep=4.8em,row sep=2.8em]
{} & \operatorname{Ker}d_{I^\bullet}^{n+1}
  \arrow[dl,dashed] \arrow[d] & {} \\
I^n \arrow[r] & I^{n+1} \arrow[r] & \cdots
\end{tikzcd}
\]

\VerdierSemanticAnchor{ch3d-3-1-4-proof-truncation}
因此，$\operatorname{Ker}d_{I^\bullet}^{n+1}$ 与
$\operatorname{Ker}d_{I^\bullet}^{n}$ 都是 $I^n$ 的直和项，因而都是
内射对象。复形
\[
\cdots\longrightarrow0\longrightarrow I^0\longrightarrow I^1
\longrightarrow\cdots\longrightarrow I^{n-1}
\longrightarrow\operatorname{Ker}d_{I^\bullet}^{n}
\longrightarrow0\longrightarrow\cdots
\]
是 $Y$ 的一个内射分解；这表明 $Y$ 的内射维数小于或等于 $n$。同理可证，
$\mathcal A$ 的每个对象 $X$ 的投射维数小于或等于 $n$。

\VerdierSemanticAnchor{ch3d-3-1-5}
\medskip\noindent\textbf{3.1.5.}\quad
设 $X$ 与 $Y$ 是 $\mathcal A$ 的两个对象。命题 (3.1.4) 表明，当范畴
$\mathcal A$ 有足够多的内射对象或投射对象时，这里定义的群
$\operatorname{Ext}_{\mathcal A}^{i}(X,Y)$ 与经典定义的群一致。

\VerdierSemanticAnchor{ch3d-3-2}
\subsection*{3.2. Yoneda 构造。}

\VerdierSemanticAnchor{ch3d-3-2-1}
\noindent\textbf{3.2.1.}\quad
设 $X$ 与 $Y$ 是 $\mathcal A$ 的两个对象。以
$E_{\mathcal A}^{n}(X,Y)$（$n\geq1$）表示正合列

\VerdierSemanticAnchor{ch3d-3-2-1-sequence}
\[
S^n(X,Y)=0\longrightarrow Y\longrightarrow Z_{n-1}
\longrightarrow Z_{n-2}\longrightarrow\cdots\longrightarrow Z_0
\xrightarrow{u_0}X\longrightarrow0.
\]
所成的集合。整数 $n$ 称为正合列 $S^n(X,Y)$ 的长度。对
$E_{\mathcal A}^{n}(X,Y)$ 的每个元素 $S^n(X,Y)$，关联如下三个复形：
\endgroup
